\chapter[Chapter \thechapter]{}
\lettrine[ante=`]{T}{'ch!} t'ch!' said Mr~Pond, clicking his tongue against his denture.
\zz
Miss~Murchison looked up from her typewriter.

\zz
<Is anything the matter, Mr~Pond?>

\zz
<No, nothing,> said the head-clerk, testily. <A foolish letter from a foolish member of your sex, Miss~Murchison.>

<That's nothing new.>

Mr~Pond frowned, conceiving the tone of his subordinate's voice to be impertinent. He picked up the letter and its enclosure and took them into the inner office.

Miss~Murchison nipped swiftly across to his desk and glanced at the registered envelope which lay upon it, open. The post-mark was <Windle.>

<That's luck,> said Miss~Murchison, to herself. <Mr~Pond is a better witness than I should be. I'm glad he opened it.>

She regained her place. In a few minutes Mr~Pond emerged, smiling slightly.

Five minutes later, Miss~Murchison, who had been frowning over her shorthand note-book, rose up and came over to him.

<Can you read shorthand, Mr~Pond?>

<No,> said the head-clerk. <In my day it was not considered necessary.>

<I can't make out this outline,> said Miss~Murchison. <It looks like <give consent to,> but it may be only 'give consideration to'—there's a difference, isn't there?>

<There certainly is,> said Mr~Pond, drily.

<P'raps I'd better not risk it,> said Miss~Murchison. <It's got to go off this morning. I'd better ask him.>

Mr~Pond snorted—not for the first time—over the carelessness of the female typist.

Miss~Murchison walked briskly across the room and opened the inner door without knocking—an informality which left Mr~Pond groaning again.

Mr~Urquhart was standing up with his back to the door, doing something or other at the mantelpiece. He turned round sharply, with an exclamation of annoyance.

<I have told you before, Miss~Murchison, that I like you to knock before entering.>

<I am very sorry; I forgot.>

<Don't let it happen again. What is it?>

He did not return to his desk, but stood leaning against the mantelshelf. His sleek head, outlined against the drab-painted panelling, was a little thrown back, as though—Miss~Murchison thought—he were protecting or defying somebody.

<I could not quite make out my shorthand note of your letter to Tewke \& Peabody,> said Miss~Murchison, <and I thought it better to come and ask you.>

<I wish,> said Mr~Urquhart, fixing a stern eye upon her, <that you would take your notes clearly at the time. If I am going too fast for you, you should tell me so. It would save trouble in the end—wouldn't it?>

Miss~Murchison was reminded of a little set of rules which Lord~Peter Wimsey—half in jest and half in earnest—had once prepared for the guidance of <The Cattery.> Of Rule Seven, in particular, which ran: <Always distrust the man who looks you straight in the eyes. He wants to prevent you from seeing something. Look for it.>

She shifted her eyes under her employer's gaze.

<I'm very sorry, Mr~Urquhart. I won't let it occur again,> she muttered. There was a curious dark line at the edge of the panelling just behind the solicitor's head, as though the panel did not quite fit its frame. She had never noticed it before.

<Well, now, what is the trouble?>

Miss~Murchison asked her question, got her answer and retired. As she went, she cast a glance over the desk. The will was not there.

She went back and finished her letters. When she took them in to be signed, she seized the opportunity to look at the panelling again. There was no dark line to be seen.

Miss~Murchison left the office promptly at half-past four. She had a feeling that it would be unwise to linger about the premises. She walked briskly away through Hand Court, turned to the right along Holborn, dived to the right again through Featherstone Buildings, made a detour through Red Lion Street and debouched into Red Lion Square. Within five minutes she was at her old walk round the square, and up Princeton Street. Presently, from a safe distance, she saw Mr~Pond come out, thin, stiff and stooping, and walk down Bedford Row towards Chancery Lane Station. Before very long, Mr~Urquhart followed. He stood a moment on the threshold, glancing to left and right, then came straight across the street towards her. For a moment she thought he had seen her, and she dived hurriedly behind a van that was standing at the kerb. Under its shelter, she withdrew to the corner of the street, where there is a butcher's shop, and scanned a windowful of New Zealand lamb and chilled beef. Mr~Urquhart came nearer. His steps grew louder—then paused. Miss~Murchison glued her eyes on a round of meat marked 4-½lb. \oldmoney{0}{3}{4}. A voice said: <Good-evening, Miss~Murchison. Choosing your supper-chop?>

<Oh! Good-evening, Mr~Urquhart. Yes—I was just wishing that Providence had seen fit to provide more joints suitable for single people.>

<Yes—one gets tired of beef and mutton.>

<And pork is apt to be indigestible.>

<Just so. Well, you should cease to be single, Miss~Murchison.>

Miss~Murchison giggled.

<But this is so sudden, Mr~Urquhart.>

Mr~Urquhart flushed under his curious freckled skin.

<Good-night,> he said abruptly, and with extreme coldness.

Miss~Murchison laughed to herself as he strode off.

<Thought that would settle him. It's a great mistake to be familiar with your subordinates. They take advantage of you.>

She watched him out of sight on the far side of the Square, then returned along Princeton Street, crossed Bedford Row and re-entered the office building. The charwoman was just coming downstairs.

<Well, Mrs~Hodges, it's me again! Do you mind letting me in? I've lost a pattern of silk. I think I must have left it in my desk, or dropped it on the floor. Have you come across it?>

<No, miss, I ain't done your office yet.>

<Then I'll have a hunt round for it. I want to get up to Bourne's before half-past six. It's such a nuisance.>

<Yes, miss, and such a crowd always with the buses and things. Here you are, miss.>

She opened the door, and Miss~Murchison darted in.

<Shall I 'elp you to look for it, miss?>

<No, thank you, Mrs~Hodges, please don't bother. I don't expect it's far off.>

Mrs~Hodges took up a pail and went to fill it at a tap in the back yard. As soon as her heavy steps had ascended again to the first floor, Miss~Murchison made for the inner office.

<I must and will see what's behind that panelling.>

The houses in Bedford Row are Hogarthian in type, tall, symmetrical, with the glamour of better days upon them. The panels in Mr~Urquhart's room, though defaced by many coats of paint, were handsomely designed, and over the mantelpiece ran a festoon of flowers and fruit, rather florid for the period, with a ribbon and basket in the centre. If the panel was controlled by a concealed spring, the boss that moved it was probably to be found among this decorative work. Pulling a chair to the fireplace, Miss~Murchison ran her fingers quickly over the festoon, pushing and pressing with both hands, while keeping her ear cocked for intruders.

This kind of investigation is easy for experts, but Miss~Murchison's knowledge of secret hiding places was only culled from sensational literature; she could not find the trick of the thing. After nearly a quarter of an hour, she began to despair.

Thump—thump—thump—Mrs~Hodges was coming downstairs.

Miss~Murchison sprang away from the panelling so hastily that the chair slipped, and she had to thrust hard at the wall to save herself. She jumped down, restored the chair to its place, glanced up—and saw the panel standing wide open.

At first she thought it was a miracle, but soon realised that in slipping she had thrust sideways at the frame of the panel. A small square of woodwork had slipped away sideways, and exposed an inner panel with a keyhole in the middle.

She heard Mrs~Hodges in the outer room, but she was too excited to bother about what Mrs~Hodges might be thinking. She pushed a heavy chair across the door, so that nobody could enter without noise and difficulty. In a moment Blindfold Bill's keys were in her hand—how fortunate that she had not returned them! How fortunate, too, that Mr~Urquhart had relied on the secrecy of the panel, and had not thought it worth while to fit his cache with a patent lock!

A few moments' quick work with the keys, and the lock turned. She pulled the little door open.

Inside was a bundle of papers. Miss~Murchison ran them over—at first quickly—then again, with a puzzled face. Receipts for securities—Share certificates—Megatherium Trust—surely the names of those investments were familiar—where had she\dots?

Suddenly Miss~Murchison sat down, feeling quite faint, the bundle of papers in her hand.

She realised now what had happened to Mrs~Wrayburn's money, which Norman Urquhart had been handling under that confiding Deed of Trust, and why the matter of the will was so important. Her head whirled. She picked up a sheet of paper from the desk and began jotting down in hurried shorthand the particulars of the various transactions of which these documents were the evidence.

Somebody bumped at the door.

<Are you in here, miss?>

<Just a moment, Mrs~Hodges. I think I must have dropped it on the floor in here.>

She gave the big chair a sharp push, effectually closing the door.

She must hurry. Anyway she had got down enough to convince Lord~Peter that Mr~Urquhart's affairs needed looking into. She put the papers back into the cupboard, in the exact place from which she had taken them. The will was there, too, she noticed, laid on one side by itself. She peered in. There was something else, tucked away at the back. She thrust her hand in and pulled the mysterious object out. It was a white paper packet, labelled with the name of a foreign chemist. The end had been opened and tucked in again. She pulled the paper apart, and saw that the packet contained about two ounces of a fine white powder.

Next to hidden treasure and mysterious documents, nothing is more full of sensational suggestion than a packet of anonymous white powder. Miss~Murchison caught up another sheet of clean paper, tipped a thimbleful of the powder into it, replaced the packet at the back of the cupboard and re-locked the door with the skeleton key. With trembling fingers she pushed the panel back into place, taking care to shut it completely, so as to show no betraying dark line.

She rolled the chair away from the door and cried out gaily:

<I've got it, Mrs~Hodges!>

<There, now!> said Mrs~Hodges, appearing in the doorway.

<Just fancy!> said Miss~Murchison. <I was looking through my patterns when Mr~Urquhart rang, and this one must have stuck to my frock and dropped on the floor in here.>

She held up a small piece of silk triumphantly. She had torn it from the lining of her bag in the course of the afternoon—a proof, if any were needed, of her devotion to her work, for the bag was a good one.

<Dearie me,> said Mrs~Hodges. <What a good thing you found it, wasn't it, miss?>

<I nearly didn't,> said Miss~Murchison, <it was right in this dark corner. Well, I must fly to get there before the shop shuts. Good-night, Mrs~Hodges.>

But long before the accommodating Messrs. Bourne \& Hollingsworth had closed their doors, Miss~Murchison was ringing the second floor bell at 110a, Piccadilly.

\divider

She found a council in progress. There was the Hon. Freddy Arbuthnot, looking amiable, Chief-Inspector Parker, looking worried, Lord~Peter, looking somnolent, and Bunter, who, having introduced her, retired to a position on the fringe of the assembly and hovered there looking correct.

<Have you brought us news, Miss~Murchison? If so, you have come at the exact right moment to find the eagles gathered together. Mr~Arbuthnot, Chief-Inspector Parker, Miss~Murchison. Now let's all sit down and be happy together. Have you had tea? or will you absorb a spot of something?>

Miss~Murchison declined refreshment.

<H'm!> said Wimsey. <The patient refuses food. Her eyes glitter wildly. The expression is anxious. The lips are parted. The fingers fumble with the clasp of the bag. The symptoms point to an acute attack of communicativeness. Tell us the worst, Miss~Murchison.>

Miss~Murchison needed no urging. She told her adventures, and had the pleasure of holding her audience enthralled from the first word to the last. When she finally produced the screw of paper containing the white powder, the sentiments of the company expressed themselves in a round of applause, in which Bunter joined discreetly.

<Are you convinced, Charles?> asked Wimsey.

<I admit that I am heavily shaken,> said Parker. <Of course, the powder must be analysed\longdash>

<It shall, embodied caution,> said Wimsey. <Bunter, make ready the rack and thumbscrew. Bunter has been taking lessons in Marsh's test, and performs it to admiration. You know all about it too, Charles, don't you?>

<Enough for a rough test.>

<Carry on then, my children. In the meanwhile, let us sum up our findings.>

Bunter went out and Parker, who had been making entries in a note-book, cleared his throat.

<Well,> he said, <the matter stands, I take it, like this. You say that Miss~Vane is innocent, and you undertake to prove this by bringing a convincing accusation against Norman Urquhart. So far, your evidence against him is almost entirely concerned with motive, bolstered up by proofs of intent to mislead enquiry. You say that your investigations have brought the case against Urquhart to a point at which the police can, and ought to, take it up, and I am inclined to agree with you. I warn you, however, that you still have to establish evidence as to means and opportunity.>

<I know that. Tell us a new one.>

<All right, as long as you know it. Very well. Now Philip Boyes and Norman Urquhart are the only surviving relations of Mrs~Wrayburn, or Cremorna Garden, who is rich, and has money to leave. A number of years ago, Mrs~Wrayburn put all her affairs into the hands of Urquhart's father, the only member of the family with whom she remained on friendly terms. On his father's death, Norman Urquhart took over those affairs himself, and in 1920, Mrs~Wrayburn executed a Deed of Trust, giving him sole authority to handle her property. She also made a will, dividing her property unequally between her two great-nephews. Philip Boyes got all the real estate and \textsterling 50,000, while Norman Urquhart took whatever was left and was also sole executor. Norman Urquhart, when questioned about this Will, deliberately told you an untruth, saying that the bulk of the money was left to him, and even went so far as to produce a document purporting to be a draft of such a will. The pretended date of this draft is subsequent to that of the Will discovered by Miss~Climpson, but there is no doubt that the draft itself was drawn out by Urquhart, certainly within the last three years and probably within the last few days. Moreover, the fact that the actual Will, though lying in a place accessible to Urquhart, was not destroyed by him, suggests that it was not, in fact, superseded by any subsequent testamentary disposition. By the way, Wimsey, why didn't he simply take the will and destroy it? As the sole surviving heir, he would then inherit without dispute.>

<Perhaps it didn't occur to him. Or there might even be other relatives surviving. How about that uncle in Australia?>

<True. At any rate he didn't destroy it. In 1925 Mrs~Wrayburn became completely paralysed and imbecile, so that there was no possibility of her ever enquiring into the disposition of her estate or making another will.

About this time, as we know from Mr~Arbuthnot, Urquhart took the dangerous step of plunging into speculation. He made mistakes, lost money, plunged more deeply to recover himself, and was involved to a large extent in the great crash of Megatherium Trust, Ltd. He certainly lost far more than he could possibly afford, and we now find, from Miss~Murchison's discoveries—of which I must say that I should hate to have to take official notice—that he had been consistently abusing his position as Trustee and employing Mrs~Wrayburn's money for his private speculations. He deposited her holdings as security for large loans, and embarked the money thus raised in the Megatherium and other wild-cat schemes.

As long as Mrs~Wrayburn lived, he was fairly safe, for he only had to pay to her the sums necessary to keep up her house and establishment. In fact, all the household bills and so on were settled by him as her man of affairs under Power of Attorney, all salaries were paid by him, and so long as he did this, it was nobody's business to ask what he had done with the capital. But as soon as Mrs~Wrayburn died, he would have to account to the other heir, Philip Boyes, for the capital which he had misappropriated.

Now in 1929, just about the time that Philip Boyes quarrelled with Miss~Vane, Mrs~Wrayburn had a serious attack of illness and very nearly died. The danger passed, but might recur at any moment. Almost immediately afterwards we find him becoming friendly with Philip Boyes and inviting him to stay at his house. While living with Urquhart, Boyes has three attacks of illness, attributed by his doctor to gastritis, but equally consistent with arsenical poisoning. In June 1929, Philip Boyes goes away to Wales and his health improves.

While Philip Boyes is absent, Mrs~Wrayburn has another alarming attack, and Urquhart hastens up to Windle, possibly with the idea of destroying the will in case the worst happens. It does not happen, and he comes back to London, just in time to receive Boyes on his return from Wales. That night, Boyes is taken ill with symptoms similar to those of the previous spring, but much more violent. After three days he dies.

Urquhart is now perfectly safe. As residuary legatee, he will receive, at Mrs~Wrayburn's death, all the money bequeathed to Philip Boyes. That is, he will not get it, because he has already taken it and lost it, but he will no longer be called upon to produce it and his fraudulent dealings will not be exposed.

So far, the evidence as to motive is extremely cogent, and far more convincing than the evidence against Miss~Vane.

But here is your snag, Wimsey. When and how was the poison administered? We know that Miss~Vane possessed arsenic and that she could easily have given it to him without witnesses. But Urquhart's only opportunity was at the dinner he shared with Boyes, and if anything in this case is certain, it is that the poison was not administered at that dinner. Everything which Boyes ate or drank was equally eaten and drunk by Urquhart and the servants, with the single exception of the burgundy, which was preserved and analysed and found to be harmless.>

<I know,> said Wimsey, <but that is what is so suspicious. Did you ever hear of a meal hedged round with such precautions? It's not natural, Charles. There's the sherry, poured out by the maid from the original bottle, the soup, fish and casseroled chicken—so impossible to poison in one portion without poisoning the whole—the omelette, so ostentatiously prepared at the table by the hands of the victim—the wine, sealed up and marked—the remnants consumed in the kitchen—you would think the man had gone out of his way to construct a suspicion-proof meal. The wine is the final touch which makes the thing incredible. Do you tell me that at that earliest moment when everybody supposes the illness to be a natural one, and when the affectionate cousin ought to be overwhelmed with anxiety for the sick man, it is natural or believable that an innocent person's mind should fly to accusations of poisoning? If he was innocent himself, then he suspected something. If he did suspect, why didn't he tell the doctor and have the patient's secretions and so on analysed? Why should he ever have thought of protecting himself against accusation when no accusation had been made, unless he knew that an accusation would be well-founded? And then there's the business about the nurse.>

<Exactly. The nurse did have her suspicions.>

<If he knew about them, he ought to have taken steps to refute them in the proper way. But I don't think he did know about them. I was referring to what you told us today. The police have got in touch with the nurse again, Miss~Williams, and she tells them that Norman Urquhart took special pains never to be left alone with the patient, and never to give him any food or medicine, even when she herself was present. Doesn't that argue a bad conscience?>

<You won't find any lawyer or jury to believe it, Peter.>

<Yes, but look here, doesn't it strike you as funny? Listen to this, Miss~Murchison. One day the nurse was doing something or the other in the room, and she had got the medicine there on the mantelpiece. Something was said about it, and Boyes remarked, <Oh, don't bother, Nurse. Norman can give me my dope.> Does Norman say, <Right-ho, old man?> as you or I would? No! He says: <No, I'll leave it to Nurse—I might make a mess of it.> Pretty feeble, what?>

<Lots of people are nervous about looking after invalids,> said Miss~Murchison.

<Yes, but most people can pour stuff out of a bottle into a glass. Boyes wasn't \binderylatin{in extremis}—he was speaking quite rationally and all that. I say the man was deliberately protecting himself.>

<Possibly,> said Parker, <but after all, old man, when \emph{did} he administer the poison?>

<Probably not at the dinner at all,> said Miss~Murchison. <As you say, the precautions seem rather obvious. They may have been intended to make people concentrate on the dinner and forget other possibilities. Did he have a whisky when he arrived or before he went out or anything?>

<Alas, he did not. Bunter has been cultivating Hannah Westlock almost to breach of promise point, and she says that she opened the door to Boyes on his arrival, that he went straight upstairs to his room, that Urquhart was out at the time and only came in a quarter of an hour before dinner-time, and that the two men met for the first time over the famous glass of sherry in the library. The folding-doors between the library and dining-room were open and Hannah was buzzing round the whole time laying the table, and she is sure that Boyes had the sherry and nothing but the sherry.>

<Not so much as a digestive tablet?>

<Nothing.>

<How about after dinner?>

<When they had finished the omelette, Urquhart said something about coffee. Boyes looked at his watch and said, <No time, old chap, I've got to be getting along to Doughty Street.> Urquhart said he would ring up a taxi, and went out to do so. Boyes folded up his napkin, got up and went into the hall. Hannah followed and helped him on with his coat. The taxi arrived. Boyes got in and off he went without seeing Urquhart again.>

<It seems to me,> said Miss~Murchison, <that Hannah is an exceedingly important witness for Mr~Urquhart's defence. You don't think—I hardly like to suggest it—but you don't think that Bunter is allowing his feelings to overcome his judgment?>

<He says,> replied Lord~Peter, <that he believes Hannah to be a sincerely religious woman. He has sat beside her in chapel and shared her hymn-book.>

<But that may be the merest hypocrisy,> said Miss~Murchison, rather warmly, for she was militantly rationalist. <I don't trust these unctuous people.>

<I didn't offer that as a proof of Hannah's virtue,> said Wimsey, <but of Bunter's unsusceptibility.>

<But he looks like a deacon himself.>

<You've never seen Bunter off duty,> said Lord~Peter, darkly. <I have, and I can assure you that a hymn-book would be about as softening to his heart as neat whisky to an Anglo-Indian liver. No; if Bunter says Hannah is honest, then she \emph{is} honest.>

<Then that definitely cuts out the drinks and the dinner,> said Miss~Murchison, unconvinced, but willing to be open-minded. <How about the water-bottle in the bedroom?>

<The devil!> cried Wimsey. <That's one up to you, Miss~Murchison. We didn't think of that. The water-bottle—yes—a perfectly fruity idea. You recollect, Charles, that in the Bravo case, it was suggested that a disgruntled servant had put tartar emetic in the water-bottle. Oh, Bunter—here you are! Next time you hold Hannah's hand, will you ask her whether Mr~Boyes drank any water from his bedroom water-bottle before dinner?>

<Pardon me, my lord, the possibility had already presented itself to my mind.>

<It had?>

<Yes, my lord.>

<Do you never overlook anything, Bunter?>

<I endeavour to give satisfaction, my lord.>

<Well then, don't talk like Jeeves. It irritates me. What about the water-bottle?>

<I was about to observe, my lord, when this lady arrived, that I had elicited a somewhat peculiar circumstance relating to the water-bottle.>

<Now we're getting somewhere,> said Parker, flattening out a new page of his note-book.

<I would not go so far as to say that, sir. Hannah informed me that she showed Mr~Boyes into his bedroom on his arrival and withdrew, as it was her place to do. She had scarcely reached the head of the staircase, when Mr~Boyes put his head out of the door and recalled her. He then asked her to fill his water-bottle. She was considerably astonished at this request, since she had a perfect recollection of having previously filled it when she put the room in order.>

<Could he have emptied it himself?> asked Parker, eagerly.

<Not into his interior, sir—there had not been time. Nor had the drinking-glass been utilised. Moreover, the bottle was not merely empty, but dry inside. Hannah apologised for the neglect, and immediately rinsed out the bottle and filled it from the tap.>

<Curious,> said Parker. <But it's quite likely she never filled it at all.>

<Pardon me, sir. Hannah was so much surprised by the episode that she mentioned it to Mrs~Pettican, the cook, who said that she distinctly recollected seeing her fill the bottle that morning.>

<Well, then,> said Parker, <Urquhart or somebody must have emptied it and dried it out. Now, why? What would one naturally do if one found one's water-bottle empty?>

<Ring the bell,> said Wimsey, promptly.

<Or shout for help,> added Parker.

<Or,> said Miss~Murchison, <if one wasn't accustomed to be waited on, one might use the water from the bedroom jug.>

<Ah!\dots of course Boyes was used to a more or less Bohemian life.>

<But surely,> said Wimsey, <that's idiotically roundabout. It would be much simpler just to poison the water in the bottle. Why direct attention to the thing by making it more difficult? Besides, you couldn't count on the victim's using the jug-water—and, as a matter of fact, he didn't.>

<And he \emph{was} poisoned,< said Miss~Murchison, >so the poison wasn't either in the jug or the bottle.>

<No—I'm afraid there's nothing to be got out of the jug and bottle department. Hollow, hollow, hollow all delight, Tennyson.>

<All the same,> said Parker, <that incident convinces me. It's too complete, somehow. Wimsey's right; it's not natural for a defence to be so perfect.>

<My God,> said Wimsey, <we have convinced Charles Parker. Nothing more is needed. He is more adamantine than any jury.>

<Yes,> said Parker, modestly, <but I'm more logical, I think. And I'm not being flustered by the Attorney-General. I should feel happier with a little evidence of a more objective kind.>

<You would. You want some real arsenic. Well, Bunter, what about it?>

<The apparatus is quite ready, my lord.>

<Very good. Let us go and see if we can give Mr~Parker what he wants. Lead and we follow.>

In a small apartment usually devoted to Bunter's photographic work, and furnished with a sink, a bench, and a bunsen burner, stood the apparatus necessary for making a Marsh's test of arsenic. The distilled water was already bubbling gently in the flask, and Bunter lifted the little glass tube which lay across the flame of the burner.

<You will perceive, my lord,> he observed, <that the apparatus is free from contamination.>

<I see nothing at all,> said Freddy.

<That, as Sherlock Holmes would say, is what you may expect to see when there is nothing there,> said Wimsey, kindly. <Charles, you will pass the water and the flask and the tube, old Uncle Tom Cobley and all as being arsenic-free.>

<I will.>

<Wilt thou love, cherish, and keep her, in sickness or in health—sorry! turned over two pages at once. Where's that powder? Miss~Murchison, you identify this sealed envelope as being the one you brought from the office, complete with mysterious white powder from Mr~Urquhart's secret hoard?>

<I do.>

<Kiss the Book. Thank you. Now then\longdash>

<Wait a sec,> said Parker, <you haven't tested the envelope separately.>

<That's true. There's always a snag somewhere. I suppose, Miss~Murchison, you haven't such a thing as another office envelope about you?>

Miss~Murchison blushed, and fumbled in her handbag.

<Well—there's a little note I scribbled this afternoon to a friend\longdash>

<\emph{In} your employer's time, \emph{on} your employer's paper,> said Wimsey. <Oh, how right Diogenes was when he took his lantern to look for an honest typist! Never mind. Let's have it. Who wills the end, wills the means.>

Miss~Murchison extracted the envelope and freed it from the enclosure. Bunter, receiving it respectfully on a developing dish, cut it into small pieces which he dropped into the flask. The water bubbled brightly, but the little tube still remained stainless from end to end.

<Does something begin to happen soon?> enquired Mr~Arbuthnot. <Because I feel this show's a bit lackin' in pep, what?>

<If you don't sit still I shall take you out,> retorted Wimsey. <Carry on, Bunter. We'll pass the envelope.>

Bunter accordingly opened the second envelope, and delicately dropped the white powder into the wide mouth of the flask. All five heads bent eagerly over the apparatus. And presently, definitely, magically, a thin silver stain began to form in the tube where the flame impinged upon it. Second by second it spread and darkened to a deep brownish-black ring with a shining metallic centre.

<Oh, lovely, lovely,> said Parker, with professional delight.

<Your lamp's smoking or something,> said Freddy.

<Is that arsenic?> breathed Miss~Murchison, gently.

<I hope so,> said Wimsey, gently detaching the tube and holding it up to the light. <It's either arsenic or antimony.>

<Allow me, my lord. The addition of a small quantity of solute chlorinated lime should decide the question beyond reach of cavil.>

He performed this further test amid an anxious silence. The stain dissolved out and vanished under the bleaching solution.

<Then it is arsenic,> said Parker.

<Oh, yes,> said Wimsey, nonchalantly, <of course it is arsenic. Didn't I tell you?> His voice wavered a little with suppressed triumph.

<Is that all?> inquired Freddy, disappointed.

<Isn't it enough?> said Miss~Murchison.

<Not quite,> said Parker, <but it's a long way towards it. It proves that Urquhart has arsenic in his possession, and by making an official enquiry in France, we can probably find out whether this packet was already in his possession last June. I notice, by the way, that it is ordinary white arsenious acid, without any mixture of charcoal or indigo, which agrees with what was found at the post-mortem. That's satisfactory, but it would be even more satisfactory if we could provide an opportunity for Urquhart to have administered it. So far, all we have done is to demonstrate clearly that he couldn't have given it to Boyes either before, during or after dinner, during the period required for the symptoms to develop. I agree that an impossibility so bolstered up by testimony is suspicious in itself, but, to convince a jury, I should prefer something better than a \binderylatin{credo quia impossibile}.>

<Riddle-me-right, and riddle-me-ree,> said Wimsey, imperturbably. <We've overlooked something, that's all. Probably something quite obvious. Give me the statutory dressing-gown and ounce of shag, and I will undertake to dispose of this little difficulty for you in a brace of shakes. In the meantime, you will no doubt take steps to secure, in an official and laborious manner, the evidence which our kind friends here have already so ably gathered in by unconventional methods, and will stand by to arrest the right man when the time comes?>

<I will,> said Parker, <gladly. Apart from all personal considerations, I'd far rather see that oily-haired fellow in the dock than any woman, and if the Force has made a mistake, the sooner it's put right the better for all concerned.>

\divider

Wimsey sat late that night in the black-and-primrose library, with the tall folios looking down at him. They represented the world's accumulated hoard of mellow wisdom and poetical beauty, to say nothing of thousands of pounds in cash. But all these counsellors sat mute upon their shelves. Strewn on tables and chairs lay the bright scarlet volumes of the Notable British Trials—Palmer, Pritchard, Maybrick, Seddon, Armstrong, Madeleine Smith—the great practitioners in arsenic—huddled together with the chief authorities on Forensic Medicine and Toxicology.

The theatre-going crowds surged home in saloon and taxi, the lights shone over the empty width of Piccadilly, the heavy night-lorries rumbled slow and seldom over the black tarmac, the long night waned and the reluctant winter dawn struggled wanly over the piled roofs of London. Bunter, silent and anxious, sat in his kitchen, brewing coffee on the stove and reading the same page of the <British Journal of Photography> over and over again.

At half-past eight the library bell rang.

<My lord?>

<My bath, Bunter.>

<Very good, my lord.>

<And some coffee.>

<Immediately, my lord.>

<And put back all the books except these.>

<Yes, my lord.>

<I know now how it was done.>

<Indeed, my lord? Permit me to offer my respectful congratulations.>

<I've still got to prove it.>

<A secondary consideration, my lord.>

Wimsey yawned. When Bunter returned a minute or two later with the coffee, he was asleep.

Bunter put the books quietly away, and looked with some curiosity at the chosen few left open on the table. They were: \workname{The Trial of Florence Maybrick}; Dixon Mann's \workname{Forensic Medicine and Toxicology}; a book with a German title which Bunter could not read; and A\@.~E\@. Housman's \workname{A Shropshire Lad.}

Bunter studied these for a few moments, and then slapped his thigh softly.

<Why, of course!> he said under his breath, <why, what a mutton-headed set of chumps we've all been!> He touched his master lightly on the shoulder,

<Your coffee, my lord.>